\documentclass[12pt]{article}
\usepackage[utf8]{inputenc}
\usepackage{graphicx}
\usepackage{hyperref}
\usepackage{listings}
\usepackage{xcolor}
\usepackage{geometry}
\usepackage{titlesec}
\usepackage{fancyhdr}

% Page Geometry
\geometry{a4paper, margin=1in}

% Code snippet styling
\definecolor{codegreen}{rgb}{0,0.6,0}
\definecolor{codegray}{rgb}{0.5,0.5,0.5}
\definecolor{codepurple}{rgb}{0.58,0,0.82}
\definecolor{backcolour}{rgb}{0.95,0.95,0.92}

\lstset{
    backgroundcolor=\color{backcolour},   
    commentstyle=\color{codegreen},
    keywordstyle=\color{magenta},
    numberstyle=\tiny\color{codegray},
    stringstyle=\color{codepurple},
    basicstyle=\ttfamily\footnotesize,
    breakatwhitespace=false,         
    breaklines=true,                 
    captionpos=b,                    
    keepspaces=true,                 
    numbers=left,                    
    numbersep=5pt,                  
    showspaces=false,                
    showstringspaces=false,
    showtabs=false,                  
    tabsize=2
}

\begin{document}

% --- Title Page ---
\begin{titlepage}
    \centering
    \vspace*{4cm}
    
    {\Huge \textbf{Scientific Calculator DevOps Pipeline}} \\
    \vspace{2cm}
    
    {\large \textbf{Abstract}} \\
    \vspace{0.5cm}
    \begin{minipage}{0.8\textwidth}
        \centering
        \textit{This project demonstrates a complete end-to-end CI/CD lifecycle for a Python-based scientific calculator. By integrating Jenkins for orchestration, Docker for containerization, and Ansible for automated deployment, the pipeline ensures that every code change is rigorously tested and seamlessly delivered to a production-ready environment.}
    \end{minipage}
    
    \vfill
    
    \begin{flushright}
        \large
        \textbf{Project By:} \\
        [Your Name] \\
        [Your Roll Number]
    \end{flushright}
\end{titlepage}

\newpage
\tableofcontents
\newpage

% --- Section 1: Introduction ---
\section{Introduction to DevOps}
DevOps is a cultural and technical movement that bridges the historical gap between Software Development (Dev) and IT Operations (Ops). Historically, these two departments functioned as silos, leading to friction during the "hand-over" phase of a project. DevOps solves this by introducing automation, shared responsibility, and continuous feedback loops.

\subsection{Core Principles}
\begin{itemize}
    \item \textbf{Continuous Integration (CI):} Developers merge their changes back to the main branch as often as possible. These changes are validated by creating a build and running automated tests against it.
    \item \textbf{Continuous Delivery (CD):} Beyond testing, the application is automatically prepared for a release to a staging or production environment.
    \item \textbf{Infrastructure as Code (IaC):} Managing and provisioning infrastructure through machine-readable definition files, rather than physical hardware configuration or interactive configuration tools.
\end{itemize}

% --- Section 2: Tech Stack ---
\section{Tech Stack \& Tools}
The choice of tools for this project was driven by industry standards and the need for a lightweight, scalable architecture.

\begin{itemize}
    \item \textbf{Git/GitHub:} Chosen for its distributed version control capabilities, allowing for local development and remote collaboration.
    \item \textbf{Jenkins:} Acting as the central automation server (Orchestrator), Jenkins handles the logic of moving code through various stages.
    \item \textbf{Docker:} Provides OS-level virtualization to deliver software in packages called containers. Containers are isolated from one another and bundle their own software, libraries, and configuration files.
    \item \textbf{Ansible:} An agentless configuration management tool that uses SSH/Local connections to automate the deployment of the Docker container.
\end{itemize}

% --- Section 3: Git Phase ---
\section{Phase 1: Source Control with Git}
Source control is the foundation of DevOps. It allows multiple developers to work on the same codebase simultaneously without overwriting each other's work.

\subsection{Implementation}
For this project, we utilized GitHub as our remote repository. We followed the \textit{Mainline Development} strategy, where the \texttt{main} branch always represents a "ready-to-deploy" state.
\vspace{1cm}
\begin{figure}[h]
    \centering
    \fbox{\begin{minipage}{0.8\textwidth} \centering \vspace{3cm} [SCREENSHOT: Git Status/Push from Terminal] \vspace{3cm} \end{minipage}}
    \caption{Git Workflow Verification}
\end{figure}

% --- Section 4: Jenkins Phase ---
\section{Phase 2: CI Orchestration with Jenkins}
Jenkins serves as the primary engine for the pipeline. It listens for events from GitHub and triggers the sequence of stages defined in the \texttt{Jenkinsfile}.

\subsection{Configuration Details}
\begin{itemize}
    \item \textbf{Webhooks:} We used \texttt{ngrok} to expose our local Jenkins instance to the internet, allowing GitHub to send "Push Notifications" via a URL.
    \item \textbf{Credentials:} Sensitive data like Docker Hub tokens were stored in Jenkins' encrypted credential store, ensuring that passwords never appear in plain text within the logs.
\end{itemize}

\vspace{1cm}
\begin{figure}[h]
    \centering
    \fbox{\begin{minipage}{0.8\textwidth} \centering \vspace{3cm} [SCREENSHOT: Jenkins Webhook and Credential Setup] \vspace{3cm} \end{minipage}}
    \caption{Jenkins Global Configuration}
\end{figure}

% --- Section 5: Jenkinsfile ---
\section{Phase 3: Pipeline as Code (Jenkinsfile)}
A key DevOps practice is "Pipeline as Code." By defining our build process in a text file (\texttt{Jenkinsfile}), we can version-control our deployment logic alongside our application code.

\subsection{Pipeline Stages}
\begin{lstlisting}[language=Groovy, caption=Declarative Jenkinsfile Structure]
pipeline {
    agent any
    environment {
        IMAGE_NAME = 'scientific-calculator'
        DOCKER_CREDENTIALS_ID = 'DockerHubCred'
    }
    stages {
        stage('Unit Testing') { ... }
        stage('Build Docker Image') { ... }
        stage('Push to Docker Hub') { ... }
        stage('Run Ansible Playbook') { ... }
    }
}
\end{lstlisting}

\vspace{1cm}
\begin{figure}[h]
    \centering
    \fbox{\begin{minipage}{0.8\textwidth} \centering \vspace{3cm} [SCREENSHOT: Jenkins Pipeline Graph View] \vspace{3cm} \end{minipage}}
    \caption{Successful Pipeline Execution}
\end{figure}

% --- Section 6: Docker ---
\section{Phase 4: Containerization with Docker}
Docker solves the "Dependency Hell" problem. By packaging the Python runtime and the \texttt{Calculator.py} script together, we ensure that the application behaves exactly the same on the Jenkins server as it does on the end-user's machine.

\subsection{Dockerfile Logic}
We utilized the \texttt{python:3.9-slim} base image to reduce the attack surface and image size. Each instruction in the Dockerfile creates a "layer," which Docker caches to speed up subsequent builds.

\vspace{1cm}
\begin{figure}[h]
    \centering
    \fbox{\begin{minipage}{0.8\textwidth} \centering \vspace{3cm} [SCREENSHOT: Docker Hub Repository with Build Tags] \vspace{3cm} \end{minipage}}
    \caption{Docker Image Management}
\end{figure}

% --- Section 7: Ansible ---
\section{Phase 5: Deployment with Ansible}
While Jenkins builds the image, Ansible is responsible for the "Last Mile" of delivery. It ensures that the host machine is in the desired state: pulling the latest image and running the container.

\subsection{Idempotency and Efficiency}
Ansible playbooks are \textit{idempotent}, meaning you can run them multiple times without changing the outcome if the system is already in the correct state. This project uses the \texttt{docker\_container} module to manage the lifecycle of the calculator application.

\vspace{1cm}
\begin{figure}[h]
    \centering
    \fbox{\begin{minipage}{0.8\textwidth} \centering \vspace{3cm} [SCREENSHOT: Ansible Execution and Running Container] \vspace{3cm} \end{minipage}}
    \caption{Automated Deployment Status}
\end{figure}

% --- Section 8: Conclusion ---
\section{Conclusion \& Future Scope}
This project successfully establishes a robust, automated pipeline for a scientific calculator. By eliminating manual steps, we have significantly reduced the risk of deployment errors.

\subsection{Future Enhancements}
\begin{itemize}
    \item \textbf{Static Code Analysis:} Integrating tools like Pylint or SonarQube to catch code quality issues before the build starts.
    \item \textbf{Orchestration:} Transitioning from Docker containers to Kubernetes pods for better scaling and self-healing capabilities.
    \item \textbf{Security Scanning:} Implementing Trivy or Snyk to scan Docker images for vulnerabilities during the push stage.
\end{itemize}

\end{document}
